\subsection{As-is domain and causal analysis.}

\subsubsection*{Reported and observed as-is practices}
ASCEND is currently associated to a Rocketry Team and an Avionics Team of the RockÉTS student club.
The Rocketry Team operates on requirements given by Rocketry Competition Organizers.
The requirements are stated in terms of the target maximum height or maximum acceleration of the rocket.
One competition imposed the requirement of launching and returning to ground an egg without breaking it, which can only be validated via demonstration.
The Team Captains convert the requirements of a competition to atomic tasks.
Tasks are performed by any club member and approved by senior members.
The club performs the PDR, CDR, and FRR review gates.

There is a large emphasis on technology reuse at the club, so most tasks are formulated as atomic edits to existing designs.
Furthermore, most of the tasks for a particular competition involve development activities like PCB manufacturing, soldering, and assembly.
One observed task involved adding a COTS to a PCB, which was specified in the task title. The task was broken down into subtasks, one of which required to connect a capacitor to the COTS, possibly prescribed by its datasheet.

The competition requirements include the radio communications technologies allowed for rocket launch commands and telemetry.
Thus, radio transceivers are some of the first COTS chosen for rocket designs.
The egg mission also introduced a payload consisting of 1U unit with a mechanically damping material to protect the egg, and heat pumps to maintain its temperature within a range inspired by an incubator machine.

Multiple lithium-ion cells were used to power the latest revision of the egg incubator payload.
To size them, a power budget was generated by consulting the datasheets of the already chosen components and extracting "on" and "off" power consumption values, then estimating the duty cycle of each component for one launch.
The duty cycle of telemetry components is estimated from the telemetry frequency requirements of the competition. For the rest of the components, it is conservatively assumed that they are always on.

The power source of the payload is verified experimentally.
The payload is fully assembled and set to provide telemetry at a realistic frequency while powered by the battery.
It is then inserted into a temperature controlled environment like a refrigerator so the heat pumps can be observed under representative operation conditions, and left to operate until the battery power runs out.
If the time from experiment start to power depletion is greater than the mission period, the power requirements are taken to be verified.
In practice, the battery life is several times larger than the mission period.

The demonstratable requirement to protect an egg has been related to the measurable variable of acceleration in the egg chamber. An IMU was placed inside this chamber, which is mechanically damped, and another in a rigid point of the payload chassis.
The difference in acceleration between the two IMUs during realistic launches is analyzed to iterative on the damping design.
Launches are recreated by manually dropping the payload from a representative height.
The drop data, and naturally an inspection of the egg, is used to characterize the damping as "too soft" or "too hard".

Maximum rocket height and maximum rocket acceleration requirements are verified at the design phase using the OpenRocket software.
This software works with abstract representations of rocket components (nose, stages, etc.) can be parametrized, combined, and then simulated.
Rocket height and acceleration requirements are taken to be verified based on the results of the simulatins, and only then does development start.

The club reports that it takes about 2 months to onboard new members to the point where they can complete tasks.
Separately, ASCEND researchers have spent 9 months on one iteration of an EPS design following a more traditional systems engineering approach, starting from requirements and currently advenced up to high-level design, with a power budget completed and battery and panel options compared but no component selections yet.

The RockÉTS club did not report following any systems engineering method.
In summary, they value early verification of rocket-specific requirements like maximum height via simulation tools, and peform multiple design iterations before developing a design. No rocket fabrication iterations were reported.
When it comes to payloads, verification comes after development even if it requires multiple fabrication iterations, possibly because the cost of development is deemed to be low.
The ASCEND research group intends to partner with RockÉTS for the CUBICS mission
in order to leverage its elecronics design (including power) and fabrication capabilities.
However, the economic, time, and risk budgets of the CUBICS mission cannot tolerate a large amount of fabrication iterations as part of the verification process. 

\subsubsection{As-is domain}

The stakeholders are the members of the ASCEND research team and RockÉTS involved in the EPS design process. The stakeholder roles identified are: Project Manager, Systems Engineer, EPS Engineer, Non-EPS Engineer, Testing Facility Technician. The Non-EPS Engineer is concerned with any other mission segment or subsystem apart from the EPS. This team structure is visible in the literature \cite{review-0529, review-1277, review-1155}, often working concurrently, and is generalized by \cite{review-0289} as "all systems architects and thematic experts".
Importantly, \cite{review-1438} acknowledges that in a small satellite project a single person may cover multiple roles, working under tight time and budget constraints. For example, at RockÉTS the EPS Engineer and the Testing Facility Technician are often the same person, and the Testing Facilities are currently limited to a temperature-controlled environment (a refrigerator).
% 0289 potential PA candidate, V&V example

\subsubsection{Stakeholder concerns}
\label{sec:stakeholder-concerns}

% Problem statement from intro.tex
Evidence fragmentation resulting in low reproducibility of analyses and low capability to detect EPS design issues early is the initial problem that motivated this project.
However, the methodology OOSEM
proposes stakeholder analysis as an examination of the limitations of the existing EPS engineering systems and workflows from each stakeholder's perspective.
The concerns below incorporate feedback from ASCEND members, the reported and observed practices at RockÉTS, and findings from the literature review.
Causal analysis is used to identify contributing factors to the concerns and improvement opportunities, derive mission requirements, and establish measures of effectiveness that can quantify the value of a potential solution.

% TODO: Update concern source locators and report-audit inputs after compression.
% Preserve historical review evidence and record new publication checks for the revised PDF.

The root concerns are listed below by stakeholder. Appendix~\ref{appendix:stakeholder-concerns} gives the full hierarchy and contributing causes. Effort means person-time.

\ConcernRole{The Project Manager is concerned with:}
% Baseline solucion.tex:31 (depth 0)
\ConcernSummary{C-PM-01}{the high effort and cost to maintain verifiable requirements;}
% Baseline solucion.tex:34 (depth 0)
\ConcernSummary{C-PM-02}{the high effort to keep financial cost budgets and estimates up to date;}
% Baseline solucion.tex:38 (depth 0)
\ConcernSummary{C-PM-03}{the high cost and effort of redesign after hardware fabrication \citep[Secs.~4--5]{handpicked-0003};}
% Baseline solucion.tex:41 (depth 0)
\ConcernSummary{C-PM-04}{the high cost of maintaining physical EPS replicas for mission troubleshooting \citep[Sec.~1]{review-0573};}
% Baseline solucion.tex:42 (depth 0)
\ConcernSummary{C-PM-05}{the high effort and cost to adapt engineering tools to the team's mission scale and resources \citep[Sec.~5.1]{buruso2026modelbased};}

\ConcernRole{The Systems Engineer is concerned with:}
% Baseline solucion.tex:48 (depth 0)
\ConcernSummary{C-SE-01}{the high effort to maintain traceability from requirements to design decisions;}
% Baseline solucion.tex:51 (depth 0)
\ConcernSummary{C-SE-02}{the high effort to communicate effectively requirement verification and validation evidence to stakeholders;}
% Baseline solucion.tex:53 (depth 0)
\ConcernSummary{C-SE-03}{the high effort to maintain EPS requirements up to date when the designs of other subsystems change;}
% Baseline solucion.tex:69 (depth 0)
\ConcernSummary{C-SE-04}{the high effort to maintain power requirements for non-EPS subsystems up to date when the EPS design changes;}
% Baseline solucion.tex:72 (depth 0)
\ConcernSummary{C-SE-05}{the low reliability, availability, and maintainability (RAM) of CubeSat designs;}
% Baseline solucion.tex:79 (depth 0)
\ConcernSummary{C-SE-06}{the low capability to detect temporary mission energy shortages before hardware integration \citep[Sec.~VI]{review-1231};}
% Baseline solucion.tex:82 (depth 0)
\ConcernSummary{C-SE-07}{(C-SE-2) the high effort to maintain consistent information across requirements, simulation models, hardware, and verification evidence \citep[Secs.~4.4 and~5.1]{buruso2026modelbased};}
% Baseline solucion.tex:87 (depth 0)
\ConcernSummary{C-SE-08}{the high effort to establish that analysis results are credible verification evidence \citep[Sec.~6]{review-1251};}
% Baseline solucion.tex:91 (depth 0)
\ConcernSummary{C-SE-09}{(C-SE-3) the high effort to keep EPS hardware interfaces and flight software architecture consistent as the design changes \citep[Secs.~3--4]{review-0853};}
% Baseline solucion.tex:96 (depth 0)
\ConcernSummary{C-SE-10}{the high effort to coordinate EPS control, telemetry collection, fault recovery, and spacecraft safe mode \citep[Sec.~V]{citation-0136};}

\ConcernRole{The EPS Engineer is concerned with:}
% Baseline solucion.tex:103 (depth 0)
\ConcernSummary{C-EPS-01}{(C-EPS-1) the high cost and effort to verify the EPS design after EPS requirements change;}
% Baseline solucion.tex:110 (depth 0)
\ConcernSummary{C-EPS-02}{the high effort to maintain traceability from EPS requirements to EPS design decisions;}
% Baseline solucion.tex:112 (depth 0)
\ConcernSummary{C-EPS-03}{the high effort to translate the verification evidence of failed gates into an updated EPS design;}
% Baseline solucion.tex:115 (depth 0)
\ConcernSummary{C-EPS-04}{the high effort to maintain EPS component requirements up to date when EPS designs change;}
% Baseline solucion.tex:130 (depth 0)
\ConcernSummary{C-EPS-05}{(C-EPS-2) the high effort to keep EPS behavioral models representative of the physical design \citep[p.~2]{handpicked-0001}\citep[Sec.~6]{review-1251};}
% Baseline solucion.tex:138 (depth 0)
\ConcernSummary{C-EPS-06}{the high uncertainty in battery energy availability and remaining lifetime under mission loads \citep[Sec.~3]{handpicked-0003}\citep[pp.~1--3]{handpicked-0001};}
% Baseline solucion.tex:142 (depth 0)
\ConcernSummary{C-EPS-07}{the low capability to detect EPS anomalies before protection thresholds are exceeded \citep[Sec.~5.1]{review-1251};}
% Baseline solucion.tex:145 (depth 0)
\ConcernSummary{C-EPS-08}{the high effort to verify that battery charging and power management decisions are safe and effective \citep[pp.~2 and~6]{handpicked-0001};}
% Baseline solucion.tex:147 (depth 0)
\ConcernSummary{C-EPS-09}{(C-EPS-3) the high effort to maintain consistent EPS command, telemetry, event, and parameter definitions across models, flight software, and ground software \citep[Secs.~3--4]{review-0853};}
% Baseline solucion.tex:150 (depth 0)
\ConcernSummary{C-EPS-10}{the high effort to specify and verify EPS command execution conditions and expected effects \citep[Secs.~4--5]{review-0853};}
% Baseline solucion.tex:152 (depth 0)
\ConcernSummary{C-EPS-11}{the high effort to communicate EPS designs during technical reviews;}

\ConcernRole{The Non-EPS Engineer is concerned with:}
% Baseline solucion.tex:158 (depth 0)
\ConcernSummary{C-NE-01}{the high cost and effort to verify a subsystem design;}
% Baseline solucion.tex:163 (depth 0)
\ConcernSummary{C-NE-02}{the high effort to integrate EPS services with flight software frameworks such as cFS and F-Prime \citep[Secs.~3--5]{review-0853}\citep[Sec.~III.B]{citation-0136};}

\ConcernRole{The Testing Facility Technician is concerned with:}
% Baseline solucion.tex:178 (depth 0)
\ConcernSummary{C-TFT-01}{the high effort to assess the feasibility and safety of a proposed verification activity that takes place in the testing facilities;}
% Baseline solucion.tex:184 (depth 0)
\ConcernSummary{C-TFT-02}{the high effort to generate verification evidence;}
